
\documentclass[submission,copyright,creativecommons]{eptcs}
\providecommand{\event}{GandALF 2025}

\usepackage{iftex}
\usepackage{breakurl}

\title{Tree Automata for Context-Sensitive Subtype Checking}

\author{Hans de Nivelle 
\institute{School of Engineering and Digital Sciences, \\
           Nazarbayev University, Astana, Kazakhstan}
\email{hans.denivelle@nu.edu.kz}
\and 
Akhmetzhan Kussainov
\institute{School of Engineering and Digital Sciences, \\
           Nazarbayev University, Astana, Kazakhstan}
\email{akhmetzhan.kussainov@nu.edu.kz}}

\def\titlerunning{Tree Automata for Context-Sensitive Type Checking}
\def\authorrunning{H. de Nivelle and A. Kussainov} 

\begin{document}

\newtheorem{theorem}{Theorem}[section]
\newtheorem{example}[theorem]{Example}
\newtheorem{definition}[theorem]{Definition}
\newtheorem{lemma}[theorem]{Lemma}
\newtheorem{question}[theorem]{Question}
\newtheorem{corollary}[theorem]{Corollary}
\newtheorem{wish}[theorem]{Wish} 
\newtheorem{algorithm}[theorem]{Algorithm}
\newtheorem{exercise}[theorem]{Exercise}

\newcommand{\proof}{{\bf Proof \\}}
\newcommand{\proofend}{{\bf End of Proof \\}}

\newcommand{\pair}{{\rm pair}}
\newcommand{\pow}{{\rm pow}}

\newcommand{\hlijn}[1]{{ \( \overline{\hspace{#1}} \) }}

\newcommand{\een}{{\bf (1) }}
\newcommand{\twee}{{\bf (2) }}
\newcommand{\drie}{{\bf (3) }}
\newcommand{\vier}{{\bf (4) }}
\newcommand{\vijf}{{\bf (5) }}

\newcommand{\Null}{{\rm Null}} 
\newcommand{\First}{{\rm First}}
\newcommand{\Trans}{{\rm Trans}} 




\newcommand{\Lucy}{{\tt LUCY}}

\newcommand{\Live}{{\rm Live}}
\newcommand{\hastype}[2]{{#1 \! \! : \! #2}}
\newcommand{\sel}[1]{?{\rm #1}}
\newcommand{\state}[1]{\langle \mbox{#1} \rangle }

\newcommand{\statrepeat}{{\bf repeat}}
\newcommand{\statblock}{{\bf block}}
\newcommand{\statbranch}{{\bf branch}}
\newcommand{\statisfalse}{{\bf isfalse}}
\newcommand{\statistrue}{{\bf istrue}}
\newcommand{\statreturn}{{\bf return}}
\newcommand{\statresolve}{{\bf resolve}}
\newcommand{\statnop}{{\bf nop}}
\newcommand{\staterase}{{\bf erase}} 


\maketitle

\begin{abstract}
   We introduce a type checking procedure for a context-sensitive type
   system. The procedure assumes that all data are trees, and
   makes use of a type of tree-automata, which can handle multisets.
   The purpose is to be able to check at compile time structs
   where the type of later fields depends on the value earlier fields.
   Such types are abundant in programming, for example if one defines
   a partial object using a boolean to determine if the object
   in the next field is defined.
   At a point where one has checked the value of earlier fields, one
   knows more about the type of later fields. This makes the
   type system context-sensitive.

   Before checking, terms are flattened. For a given
   function, all overload candidates are inserted into
   an intermediate representation.
   After that, a fixed point is computed on the intermediate
   representation, while at the same time 
   all non-viable overloads are removed.
   If everything goes well, this result in a
   checked representation, from which code can be generated.
   States are approximated by tuples of states of a tree automaton,
   and arrays are cut into slices that are treated as multisets.
   In this way, we can handle loops that change the properties
   of an array.
\end{abstract}

\section{Introduction}

This paper originates from a project to design and implement
a programming language for the implementation
of logical algorithms, as they occur for example in theorem provers
or interactive verification systems.
In programming, one often has a situation where a 
{\bf struct} (or record) type has a tree-like type definition:
The type of the first field is fixed, but the 
types of later fields depends on the values of the 
fields that occur before it. 
The simplest such case is \verb|std::optional|, which consists
of a boolean that indicates whether the value is there.
Second simplest case is \verb|std::variant<T1, ..., Tn>|, which contains
an integer indicating which type \verb|Ti| is currently present.
A third case appears in parsing, when one uses attribute grammars.
Symbols in an attribute grammar consist of a tag (specifying
which symbol is present), optionally followed by an attribute, whose
type depends on the symbol tag. For example if the symbol tag is 
string, identifier, or scan error, 
the attribute will have type \verb|std::string|. 
If the symbol tag is double, the attribute must be \verb|double|. 
The fourth case arises from implementation of logic: A propositional
formula consists of a logical operator, followed by either a string, a
single formula, two formulas, or a sequence of formulas,
dependent on the logical operator in front. 
Current static type systems do not support such type dependencies. 
After checking the first field, access to the following fields has take
place through operations that can fail at run time. 
The access of these fields can be dynamically type checked, but
it is beyond current static type systems. 

We want to design and implement type checking
procedures that can check such dependent 
types for realistic programs. In this paper, we present a design
that we believe will work in practice. 
It uses a novel type of tree-automata,
that we call \emph{deterministic multiset automata}. Their
distinguishing feature is that they are able to process multisets, 
up to a limit. The limit is a result of the finite number of states.
We use DMA-s to model array ranges.
If one uses an index variable $ i $ to change the values in an
array $ A $ with a for-loop, we will model $ A[ 0 \ldots i \rangle $ 
and $ A[ i \ldots \rangle $ as states of a DMA.

In Section~\ref{Sect_data_terms}, we introduce the data 
that our procedure works on. These are trees over a small
set of primitive types. Proper trees are constructed by a tuple
operator, and by an array operator. 
Our type system does not aim to analyze arithmetical properties
of numeric types, as done in \cite{JhaVaz2021ref}. 
DMA's can handle properties of form `$i$ is interval $ [2,5] $'
but not '$i$ is even'. Also, they don't support arithmetic,
so a DMA would not be able to determine that $ j := i + 1 $ 
implies that `$j$ will be in the interval $ [3,6]. $' 

During type checking, every function body is checked independently.
Type checking is always combined with overload resolution. 
The source code contains imprecise names, which can be mapped
to zero, one or any number of candidate implementations.
The typechecker checks which available implementations are 
guaranteed to succeed (whose preconditions are guaranteed to be met).
It removes the candidates whose preconditions are not met. 
If no implementation remains, the type checker reports an error. 
If more than one available implementation remains, the type checker
has to pick the most specific implementation. If no implementation
is most specific, the result is a type checking error (ambiguous overload).

Standard compilers perform type checking and intermediate code generation
at the same time.
Expressions are type checked, the overloads are resolved, and the 
expression is translated into intermediate code. 
(See \cite{WiSe2023})
For our type system, this is not possible, because
subtypes can change (disappear)
during loop iteration. As a consequence, we need to
use an intermediate
representation in which overloads are not resolved. Instead, 
for every imprecise statement, we
insert all overload candidates into the intermediate representation. 
Consider for example:
\begin{verbatim}
   x = 0
   while ...
      print(x)
         // There are two overloads printE (for even numbers),
         // and print (for arbitrary numbers). 
      x = ( ?succ, ( ?succ, x )).
\end{verbatim}
In order to resolve \verb| print(x) |into \verb| printE(x) |, 
one must know that all
operations on \verb|x| in the loop preserve evenness.
Some of the operations appear after the print statement in the code.
So, we will insert both versions of \verb|print| into the 
intermediate representation,
approximate all values that $ x $ can have using DMAs, and verify
that $ x $ is always even. In the end, both versions of
\verb|print| are eligible, but \verb|printE| is more specific,
so that it can be picked. 

Unfortunately, approximating values of variables is more
complicated than we sketched above. It is not sufficient
to collect all possible values of all variables at each program point
and map them into DMA states.
We need to collect combinations of possible values of variables. 
This is due to the fact that there may be dependencies between variables,
that would otherwise get lost in the process.
For example, if one compares two propositional formulas, one
typically starts by comparing the operators. If the operators are distinct,
the formulas are not equal. After that, one can do a case analysis
on the operator. If one formula has a binary operator on top,
one can safely access the left and right subformula in both formulas,
because one knows that the operators are equal. 
This information would be lost, if one would keep at each
program point only possible
values, instead of tuples of combinations of values.

In the next section, we define \emph{data terms}. 
Their definition is straightforward. They are built-up
from primitive types using a tuple constructor, and an array constructor. 

In the next section~\ref{Sect_DMA}, we define
deterministic multiset automata. 
They are a straightforward generalization of standard,
deterministic tree automata (DFTA in \cite{TATA2007})
to include multiset nodes. Evaluation is always bottom-up. 

In Section~\ref{Sect_intermediate}, we present an intermediate
representation. It uses flattened expressions, which
is quite common in compilers, see e.g.
\cite{LLVM2004}, \cite{Muchnick_CDI_1997}, or \cite{WiSe2023}.
Different from the standard approach, it uses a 
$ \statresolve $ statement, which for a given imprecise
statement, contains all candidate overloads, together with 
their preconditions, encoded as DMAs. 
The representation does not use flow graphs,
because the resolve statement is not naturally represented in
flow graphs. Our recursively defined format is slightly more
complicated than standard flow graphs, but much more flexible.
In Section~\ref{Sect_range_states}, we approximate states by 
mappings from variables to states of multiset automata. 
In addition, for every variable that can serve as an index of an array, 
a state keeps track of the values before the index and after the index
as separate multisets. The approximations are called 
\emph{range states} because they keep track of ranges in
arrays. We explain the effect of different statements
on range states. 

\section{Data Terms}
\label{Sect_data_terms}

We assume that all data are trees over a fixed set of primitive
types. Most primitive types are standard.
We use 
$ {\bf char}, {\bf u64}, {\bf bool}, {\bf integer}, $ and
$ {\bf float}. $ Type $ {\bf u64} $ consists of 64-bit integers,
which can be used to index arrays. Type $ {\bf integer} $ 
consists of unbounded precision integers.  
In addition to the standard primitive types, there is a type
called $ {\bf selector}, $ which consists of named constants.
It is similar to an enumeration type, but with two differences:
Firstly, {\bf selectors} cannot be converted to (or obtained from) integers,
and secondly, there is only one $ {\bf selector} $
type. There is no need to allow
the user to define new enumeration types, because our type
system supports subtypes. Their intended use is to represent
tags that determine the types of later fields, 
like for examples the symbols in an attribute
grammar, or logical operators.
Compound trees are built using two constructors: The first
constructor constructs trees of fixed width. Once constructed,
the width cannot change.
The second constructor constructs trees of unbounded width,
with can grow and shrink like a vector. They are initially
constructed empty, after which elements can be inserted later.
\begin{definition}
   \label{Def_primtype} 
   We assume the following primitive types:
      {\bf bool}, {\bf char}, {\bf selector},
      {\bf u64}, {\bf integer}, {\bf double}.
   We call {\bf bool}, {\bf char}, and {\bf selector}
   \emph{tag types}, because they can be used as tags
   for selecting options.
   We assume that $ {\bf u64}, \ {\bf char} $ and $ {\bf bool} $
   are ordered types with the usual order.
\end{definition}

\begin{definition}
   \label{Def_data_term}
   We define the set of \emph{data terms} 
   by recursion as follows:
   \begin{itemize}
   \item
      If $ d $ is an element of one of the primitive types
      listed in Definition~\ref{Def_primtype},
      then $ d $ is a data term.
   \item
      If $ d_1, \ldots, d_n $ with $ n \geq 0 $ are
      data terms, then
      $ ( d_1, \ldots, d_n ) $
      and
      $ [ d_1, \ldots, d_n ] $ are also data terms. 
      We call $ ( d_1, \ldots, d_n ) $ a \emph{tuple}, and
      $ [ d_1, \ldots, d_n ] $ an \emph{array}.
   \end{itemize}
\end{definition}
The difference between $ (d_1, \ldots, d_n) $ and
$ [ d_1, \ldots, d_n ] $ is that the former represents
a fixed sized tuple, while the latter represents a vector.
In a program, the elements of a tuple $ (d_1, \ldots, d_n) $ can be
accessed by field access. For each field access instruction,
it is known which field it accesses.
At the same time, an array term $ [ d_1, \ldots, d_n ] $ can be indexed
by a variable of type $ {\bf u64}. $ 
\begin{definition}
   \label{Def_width}
   For a data term $ t, $ we define $ {\bf width}(t) $
   as the maximal size of a tuple occurring in $ t. $
   If $ t $ contains no tuples, then $ {\bf width}(t) = 0. $
\end{definition}
Because a program is always finite, 
there exists, for every program, an upperbound on the size
of the tuples that occur in it.
Hence, there always exists a natural number $ n, $ s.t.
all data terms that the program operates on, have
width at most $ n. $

\begin{example}
   Assume that we want to represent logical formulas  
   using operators $ \neg, \wedge, \vee, \rightarrow, \leftrightarrow, $ 
   where $ \neg $ is unary, $ \wedge $ and $ \vee $ are $ n $-ary, 
   and $ \leftrightarrow,\rightarrow $ are binary.
   For simplicity, we assume that atoms are represented
   by arrays of characters.

   The atom $ v0 $ can be represented by
   $ ( \sel{atom}, [ 'v', '0' ] ). $ 
   It is a tuple consisting of the selector constant
   \sel{atom} and an array, containing characters 'v' and '0'.
   
   The negation of $ v0 $ is $ ( \sel{neg}, v0 ). $ 
   Formula $ v1 \rightarrow v2 $ can be represented as
   $ ( \sel{implies}, ( \sel{atom}, [ 'v', '0' ] \, ), \,
                      ( \sel{atom}, [ 'v', '1' ] )). $ 
   Formula $ v_1 \wedge v_2 $ be be represented as
   $ ( \sel{and}, [ ( \sel{atom}, [ 'v', '0' ] \, ), \, 
                    ( \sel{atom}, [ 'v', '1' ] \, ) \, ] \, ). $ 
   The second element after an \sel{and}, or \sel{or} is always an
   array term. 
\end{example}
It can be seen from the example that, in constrast to most other
programming languages, logical operators are not represented by
constructors or function symbols, but by 
{\bf selector} constants. All formulas are constructed by
the same tuple constructor. For propositional formula,
the biggest arity that a tuple
can have, is three.
$ n $-ary formulas are represented by a binary tuple, whose
second element is an array.

\section{Multiset Automata}
\label{Sect_DMA}

Preconditions and conditions in programs are represented by
deterministic multiset automata (DMAs).
A DMA has a finite set of states. 
We use a standard bottom-up semantics (see \cite{TATA2007}),
and we will assume that tree automata are \emph{complete}, i.e.
never get stuck. 
We will merge all preconditions and conditions into a single
DMA. This automaton will use different sets of accepting states
for different (pre)conditions, 
so we do not make the accepting states part of the
automaton. 

On tuples, a DMA behaves as usual: For 
$ ( t_1, \ldots, t_n ), $ once the states for the
$ t_i $ are known, the state for the
tuple can be determined. Since the possible states
for the $ t_i $ are finite in number, and the maximal
value of $ n $ is a priori known, the behaviour
of a DMA on tuples can be finitely expressed.

Arrays will be treated as multisets. In our observation,
all types and subtypes that 
we want to express, can be expressed as set properties
(either at least one element, or all elements must have some property). 
If a type property
cannot be expressed as a set, this means that order matters.
In that case, values at different positions have different meanings,
which implies that they should have been in different fields.
In order to define set transitions, it is sufficient to 
define a state for the empty set, and a binary function for
insertion of a single element. Insertion must be idempotent, 
and order insensitive. This results in two axioms. However, it turned out
that idempotence was not needed, so we omitted it.
This results in a limited form of multiset semantics. 
As a consequence, DMAs can count a bit, but 
with an upperbound, because the set of states is always finite.

\begin{definition}
   \label{Def_DMA}
   We define a \emph{deterministic multiset automaton (DMA)}
   $ {\cal A} $ as a triple
   $ ( Q, \delta, N ), $ where $ Q $ is a finite (nonempty) set of states,
   and $ \delta $ is a transition function with the following
   properties:
   \begin{itemize}
   \item
      If $ c $ is a constant of one of the primitive types
      defined in Definition~\ref{Def_primtype}, then
      $ \delta(c) \in Q. $
   \item
      If $ q_1, \ldots, q_n \in Q $ and $ 0 \leq n \leq N, $
      then
      $ \delta( q_1, \ldots, q_n ) \in Q. $
   \item
      $ \delta( ) \in Q, $
   \item
      If $ q_1,q_2 \in Q, $ then
      $ \delta( q_1; q_2 ) \in Q. $
   \end{itemize}
   The intended meaning of $ ; $ is multiset insertion, with
   the inserted element to the right.
   Technically, $ \delta $ should be viewed as distinct overloads,
   namely
   $ \delta_T $ for primitive type $ T, $ \
   $ \delta_{()} $ for tuples,
   $ \delta_{\emptyset} $ for the empty multiset, and $ \delta_{;} $
   for multiset insertion.
   W.r.t to multiset insertion,
   $ \delta $ must have the following property,
   expressing irrelevance of insertion order: \ 
   $ 
       \delta( \delta( q_1; q_2 ); q_3 ) =
       \delta( \delta( q_1; q_3 ); q_2 ).
   $ 
\end{definition}
As already said before, we are not requiring the property
$ \delta( \delta( q_1, q_2 ); q_2 ) = \delta( q_1; q_2 ), $ which
would enforce set semantics. 

We implicitly assume that $ \delta_T $ on primitive types is
easily representable. The problem is that on an infinite, primitive type,
like $ T = {\bf integer}, \ \ \delta_T $ could represent an 
undecidable property. We assume that on a primitive type $ T, $ 
$ \delta_T $ can be represented by a finite set of values or intervals.

Using a DMA $ {\cal A}, $ every term can be mapped to a state of
$ {\cal A}. $ We will call the resulting function $ Q_{\cal A}: $
\begin{definition}
   For a deterministic tree automaton $ {\cal A} = ( Q, \delta, N ) $,
   we define a function $ Q_{\cal A} $ from data terms to $ Q $
   by recursion as follows:
   \begin{itemize}
   \item
      If $ c $ is a constant of primitive type,
      then $ Q_{\cal A}(c) = \delta(c) $.
   \item
      $ Q_{\cal A}( \ ( d_1, \ldots, d_n ) \ ) =
         \delta( \, Q_{\cal A}(d_1),\ldots,Q_{\cal A}(d_n) \, ) $.
   \item
      $ Q_{\cal A}( \, [ \, ] \, ) = \delta( \, ) $.
   \item
     If $ n > 0, $ then
     $ Q_{\cal A}( \, [ d_1, \ldots, d_n ] \, ) =
          \delta( \, Q_{\cal A}( \, [ d_1, \ldots, d_{n-1} ] \, ); \,
                     Q_{\cal A}(d_n) \, ). $
   \end{itemize}
\end{definition}
Our definition of DMA does not include accepting states.
Our type checking algorithm starts by collecting all possible overloads
for each imprecise name in a translation unit.
Each overload has a type/precondition which
is defined by a DMA. We compute the product of all these automata,
where each contributing automata defines a set of accepting states
in the product automaton. As a consequene, we end up with a single
DMA that is used with different sets of accepting states.

\begin{definition}
   \label{Def_acc_cond}
   Let $ {\cal A} = ( Q, \delta, N ) $ be a DMA.
   An \emph{acceptance condition} $ F $ is a unary predicate on $ Q. $
   An \emph{acceptance relation} $ R $ is a binary predicate on $ Q. $
\end{definition}
Acceptance conditions are used for expressing conditions and
preconditions. Acceptance relations are used for expressing 
conditions only.
\begin{example}
   Propositional formulas are expressed by the following DMA using
   acceptance condition $ \{ \state{form} \}: $ 
   \[
      \begin{array}{lll}
         \delta( \sel{atom} ) & = & \state{atom} \\
         \delta( \sel{neg} ) & = & \state{neg} \\
         \delta( \sel{implies} ) = 
            \delta( \sel{equiv} ) & = & \state{binary} \\
         \delta( \sel{and} ) = 
            \delta( \sel{or} ) & = & \state{nary} \\
         \mbox{ for every {\bf char} } c : \ \ \delta( c ) & = & \state{char} \\
         \\
         \delta( \, ) & = & \state{empty} \\
         \delta( \state{empty}; \state{char} ) = 
         \delta( \state{allchar}; \state{char} ) & = & \state{allchar} \\
         \delta( \state{empty}; \state{form} ) =
         \delta( \state{allform}; \state{form} ) & = & \state{allform} \\
         \\  
         \delta( \state{atom}, \state{allchar} ) & = & \state{prop} \\
         \delta( \state{neg}, \state{prop} ) & = & \state{prop} \\
         \delta( \state{binary}, \state{prop}, \state{prop} ) & = & \state{prop} \\
         \delta( \state{nary}, \state{allprop} ) & = & \state{prop} \\
      \end{array}
   \]
   All transitions that are not listed go to a trap state $ \state{trap}. $ 
\end{example} 
In the example, there is a single trap state $ \state{trap}. $ In practice,
it is convenient to use more than one state for data terms that are not part
of a formula, so that not all information is lost when a formula
is stored in some other data structure.

\noindent
Taking the product of two automata is easy in principle, but
one has to deal with the case where the two automata have different
width:

\begin{definition}
   \label{Def_inflation}
   Let $ {\cal A} = ( Q, \delta, N ). $
   Assume that $ N' > N. $
   The result of \emph{inflating} $ {\cal A} $ to $ N' $
   is the DMA $ {\cal A}' = ( Q', \delta', N' ), $ satisfying: 
   \begin{itemize}
   \item
      $ Q' = \{ Q \} \cup \{ q' \}, $ with $ q' \not \in Q $ a new state.
   \item
      If $ n \leq N, $ then $ \delta'( q_1, \ldots, q_n ) = \delta(q_1, \ldots, q_n). $
      If $ N < n \leq N', $ then $ \delta( q_1, \ldots, q_n ) = q'. $
   \item
      If one of the $ q_i = q', $ then $ \delta(q_1, \ldots, q_n ) = q'. $ 
   \item
      If $ q_1 = q' $ or $ q_2 = q', $ then $ \delta(q_1;q_2) = q'. $ 
   \end{itemize}
\end{definition}

\begin{definition}
   \label{Def_product_automaton}
   Let $ {\cal A}_1 = ( Q_1, \delta_1, N_1) $ and
   $ {\cal A}_2 = ( Q_2, \delta_2, N_2 ) $
   be deterministic tree automata.
   We define the product automaton
   $ {\cal A}_1 \times {\cal A}_2 =
        ( \, Q_1 \times Q_2, \ \delta', \ {\bf max}(N_1,N_2) \, ) \ $
   as follows:
   If $ N_1 \not = N_2, $ then first inflate the smaller
   automata to $ \max(N_1,N_2). $
   After that, construct $ \delta' $ as follows:
   \begin{itemize}
   \item
      For a constant $c$ of primitive type,
      $ \delta'(c) = ( \, \delta_1(c), \delta_2(c) \, ). $
   \item
      $ \delta'( \, (q_1,q'_1), \ldots, (q_n,q'_n) \, ) =
      ( \, \delta_1( q_1, \ldots, q_n), \, \delta_2( q'_1, \ldots, q'_n ) \, ). $
   \item
      $ \delta'( ) = ( \, \delta_1( ), \delta_2( ) \, ). $
   \item
      $ \delta'( \, ( q_1,q'_1); ( q_2, q'_2) \, ) =
        ( \, \delta_1( q_1; q_2), \delta_2( q'_1, q'_2) \, ). $ 
    \end{itemize}
\end{definition}
The product can be straightforwardly generalized to more than two automata.
If $ {\cal A}_1 $ or $ {\cal A}_2 $ has an acceptance condition,
then this acceptance conditions gets 'inherited' by the product automaton
as follows:
\begin{definition}
   \label{Def_acc_proj}
   Let $ F $ be an acceptance condition.
   $ {\bf proj}_i $ is an acceptance condition on tuples
   of arity at least $ i, $ defined from
   $ {\bf proj}_i( \, ( q_1, \ldots, q_n ) \, ) \Leftrightarrow
     F(q_i). $
\end{definition}
Using projection, acceptance conditions can be carried over to the 
product automaton. 
This has one important consequence, namely that
all Boolean combinations of properties definable
by a DMA + acceptance condition can be decided
by constructing the product automata, constructing
a Boolean combination of projections, and verifying it for
all states. At this point, we are relying on the fact that
$ \delta_T $ has low complexity on primitive type $ T. $ 

In particular, we will need to check inclusion between subtypes, 
in cases where more than one overload candidate remains.
In that case, one overload must be strictly more specific than
all other overload candidates.

\section{Intermediate Code Representation} 
\label{Sect_intermediate}

We introduce a format for representation of intermediate code.
For our type system, intermediate code cannot be generated in the
standard way.
Standard compilers create intermediate code immediately
after type checking. In the intermediate representation, all operations
have become unambiguous. This is possible because every variable
receives its type at the moment of initialization. 
Both LLVM (\cite{LLVM2004}) and GIMPLE (the intermediate representation
used by GCC, \cite{GIMPLE}) rely on unambiguous types for every value.
Unfortunately, this is not possible for our type system, 
because subtypes in general cannot be read off from declarations.
In the presence of loops, subtypes have to be obtained by means
of a fixed point computation. 

When an operation is ambiguous, all candidate overloads have to 
be inserted into the intermediate representation, after which 
those overloads whose preconditions are not met, need to be removed during
fixed point computation. 

\begin{definition}
   \label{Def_flatterm}
   We define \emph{flatterms}:
   \begin{itemize}
   \item
      If $ v $ is a local variable, then $ v $ is a flatterm.
   \item
      If $ c $ is a constant of a primitive type,
      then $ c $ is a flatterm.
   \item
      If $ {\cal A} $ is a DMA, and $ F $ is a state
      condition of $ {\cal A}, $ then $ {\bf some}({\cal A}, F ) $ is
      a flatterm.
   \item
      If $ v_1, \ldots, v_n $, are variables, $ f $ is an identifier,
      then $ f(v_1, \ldots, v_n) $ is a flatterm.
   \item
      If $ v $ is a variable, then
      $ {\bf len}(v), \ {\bf not}(v), \ {\bf neg}(v) $ 
      are flatterms.
   \item
      If $ v_1 $ and $ v_2 $ are variables,
      then 
      $ {\bf eq}(v_1,v_2), \ {\bf ne}(v_1,v_2), \
         {\bf lt}(v_1,v_2), \ {\bf gt}(v_1,v_2), $ \\
      $ {\bf le}(v_1,v_2), \ {\bf ge}(v_1,v_2), \
        {\bf add}(v_1,v_2), \ {\bf sub}(v_1,v_2)  $
      are flatterms.
  \item
      If $ v $ is a variable,
      $ {\cal A} $ is a DMA, and $ F $ is a state condition, then
      $ {\bf check}( v, {\cal A}, F ) $ is a flatterm.
   \item
          if $ v $ and $ w $ are variables,
          then $ {\bf ind}(v,w) $ is a flatterm.
   \item
      If $ v $ is a variable, and $ f $ is an identifier,
      then $ {\bf fld}_f(v) $ is a flatterm.
   \item
      If $ v_1, \ldots, v_n $ are variables with $ n \geq 0 $,
          then $ ( v_1, \ldots, v_n ) $ and
                   $ [ v_1, \ldots, v_n ] $ are flatterms.
   \end{itemize}
\end{definition}

\noindent
The meanings are as follows: If $ v $ is an array, then $ {\bf len}(v) $
equals its size. The intuitive meaning of $ {\bf ind}(v,w) $ is $ v[w]. $
The intuitive meaning of $ {\bf fld}_f(v) $ is $ v.f. $
Flatterm $ {\bf some}( {\cal A}, F ) $ non-deterministically creates a
data term $ t, $ s.t. $ F( \, Q_{\cal A}(t) \, ), $ while
$ {\bf check}( t, {\cal A}, F ) $ returns a boolean representing
truth of $ F( \, Q_{\cal A}(t) \, ). $
\begin{definition}
   \label{Def_program_point}
   We assume a countable set $ \Pi $ of program points.
\end{definition}
It is not very important how program points are represented.
Since statements defined below are trees, one can represent them as
paths.
\begin{definition}
   \label{Def_statement}
   We recursively define \emph{statements}:
   For a given statement, we assume that every substatement
   has a unique program point.
   A statement has one of the following forms:
   \begin{itemize}
   \item
      A scalar assignment of $ v \leftarrow t, $ in which $ v $ is a variable
      and $ t $ is a flatterm.
   \item
      A array assignment of $ v[w] := t $, in which
      $ v $ and $ w $ are variables,
      and $ t $ is a flatterm.
      Note that array assignment is a single operation
      with three arguments.
   \item
      If $ s_1, \ldots, s_n $ are statements, then
      $ \statblock(s_1, \ldots, s_n), \
        \statrepeat(s_1, \ldots, s_n), $ and 
        $ \statbranch(s_1, \ldots, s_n) $ are also statements.
   \item
      If $ v $ is a variable and
      $ s $ is a statement, then
      $ \statistrue \, v $ and $ \statisfalse \, v $ are also statements.
   \item
      $ \statnop $ is a statement.
  \item
      If $ v $ is a variable, then $ {\bf return} \, v  $ is a statement.
   \item
      $ v. {\rm push\_back}( w ) $
      with $ v,w $ variables is a statement.
   \item
      $ v. {\rm pop\_back}( ) $ is a statement.
   \item
      A field assignment $ w_1. f := w_2, $
      where $ f $ is a scalar field name,
      and $ w_1,w_2 $ are variables.
   \item
      If $ \overline{v} $ is a sequence of variables,
      and for each $ i $ with $ 1 \leq i \leq n, $ \ \
      $ {\cal A}_i $ is a sequence of automata,
      $ F_i $ is a sequence of state conditions,
      s.t. $ \| v \| = \| {\cal A}_i \| = \| F_i \|, $
      each $ s_i $ is a statement, then \\
      $ \statresolve( \overline{v}, ( {\cal A}_1, F_1 ) / s_1, \ldots,
                                    ( {\cal A}_n, F_n ) / s_n ) \, ) $
      is also a statement.
   \item
      If $ v $ is a variable, then $ \staterase \, v $ is a statement.
   \end{itemize}
\end{definition}

\noindent
Both $ \statrepeat $ and $ \statbranch $ are non-deterministic.
The $ \statrepeat $ statement non-deterministically skips
$ s_1, \ldots, s_n $ or starts at $ s_1. $
Similarly, $ \statbranch $ non-deterministically executes one of its
$ s_i. $ We chose for non-determinism because analyzing
non-deterministic code is not harder than deterministic code,
due to the fact that overapproximation will introduce non-determinism anyway.
In order to model deterministic code, one has to put
a $ \statistrue $ or $ \statisfalse $ at the beginning of every branch,
at the beginning of very \statrepeat, and immediately after every
\statrepeat.
Both statements quietly fail when the condition is not met.

In the $ \statresolve $ statement, the variables $ \overline{v} $ are used
for type checking. Overload $ ( {\cal A}_i, F_i ) / s_i $ is viable,
if for each variable in $ \overline{v}, $ the value, when
run with the corresponding automaton in $ {\cal A}_i, $ reaches a state
that satisfies the state corresponding condition.
\noindent
The difference between {\bf abort} and {\bf error} is
that $ {\bf error} $ is checked during type checking,
while $ {\bf abort} $ is checked at run time.

\begin{definition}
   \label{Def_state}
   A \emph{state} $ \sigma $ is a finite mapping from variables to
   data terms.
\end{definition}

We now define a function $ {\bf eval}_{\sigma}(t) $
that evaluates a flatterm $ t $ in state $ \sigma. $
Because the semantics of flatterms is non-deterministic,
$ {\bf eval}_{\sigma}(t) $ defines a set of possible values.

\begin{definition}
   \label{Def_semantics}
   For a given statement $ s, $ we assume two functions
   $ \Phi_{\bf in} $ and $ \Phi_{\bf out} $ that map
   program points to set of states.
   Both functions are defined for the program points occurring
   in $ s. $ They are defined together as the least fixed point
   satisfying the following conditions:
\begin{itemize}
\item
   If $ \pi $ contains a $ \statnop, $ then
   $ \Phi_{\bf in}( \pi ) \subseteq \Phi_{\bf out}( \pi ). $
\item
   If $ \pi $ contains $ \statblock( s_1, \ldots, s_n), $
   then let $ \pi_1, \ldots, \pi_n $ be the positions of
   $ s_1, \ldots, s_n. $
   We have $ \Phi_{\bf in}( \pi ) \subseteq \Phi_{\bf in}( \pi_1). $
   For every $ i $ with $ 1 \leq i \leq n-1, $ we have
   $ \Phi_{\bf out}( \pi_i ) \subseteq \Phi_{\bf in}( \pi_{i+1} ), $ and
   $ \Phi_{\bf out}(\pi_n) \subseteq \Phi_{\bf out}(\pi). $
\item
   If $ \pi $ contains $ \statrepeat( s_1, \ldots, s_n ), $ then
   let $ \pi_1, \ldots, \pi_n $ be the positions of
   $ s_1, \ldots, s_n. $
   We have $ \Phi_{\bf in}( \pi ) \subseteq \Phi_{\bf in}( \pi_1). $
   For every $ i $ with $ 1 \leq i \leq n-1, $ we have
   $ \Phi_{\bf out}( \pi_i ) \subseteq \Phi_{\bf in}( \pi_{i+1} ). $
   In addition we have $ \Phi_{\bf in}(\pi) \subseteq \Phi_{\bf out}(\pi), $ and
   $ \Phi_{\bf out}( \pi_n ) \subseteq \Phi_{\bf in}(\pi). $
\item
   If $ \pi $ contains $ \statbranch( s_1, \ldots, s_n ), $
   then let $ \pi_1, \ldots, \pi_n $ be the positions of
   $ s_1, \ldots, s_n. $
   For every $ i, $ we have
   $ \Phi_{\bf in}( \pi ) \subseteq \Phi_{\bf in}( \pi_i ), $ and
   $ \Phi_{\bf out}( \pi_i ) \subseteq \Phi_{\bf out}(\pi). $
\item
   If $ \pi $ contains $ \statistrue \, v, $ then
   $ \Phi_{\bf in}( \pi ) \cap
     \{ \sigma \mid \sigma(v) = {\bf t} \} \subseteq
     \Phi_{\bf out}( \pi). $
\item
   $ \statisfalse $ is analogous to $ \statistrue. $
\item
   If $ \pi $ contains a statement of form $ v \leftarrow t $ with
   $ t $ a flatterm, then
   \[ \{ \sigma \cup \{ v := d \} \, \} \mid
           \sigma \in \Phi_{\bf in}(\pi) \mbox{ and }
           d \in {\bf eval}_{\sigma}(t) \, \} \subseteq
           \Phi_{\bf out}( \pi ). \]
   Function $ {\bf eval}_{\sigma}(t) $ is a function
   that (nondeterministically) evaluates flatterms. 
\item
   If $ \pi $ contains $ \statresolve( c_1/s_1, \ldots, c_n/s_n ), $
   then let $ \pi_1, \ldots, \pi_n $ be the positions of
   $ s_1, \ldots, s_n. $ Similary, let $ \rho_1, \ldots, \rho_n $
   be the positions of $ c_1, \ldots, c_n. $
   We have $ \Phi_{\bf in}(\pi) \subseteq \Phi_{\bf in}( \pi_i ) $ and
   $ \Phi_{\bf in}( \pi ) \subseteq \Phi( \rho_i ). $

   We also have
   \[ \bigcap \Phi_{\bf out}( \pi_i ) \subseteq \Phi_{\bf out}( \pi ). \]
\end{itemize}
\end{definition}
The last condition is a bit tricky.
It is possible that some of the $ s_i $ use temporary
variables that other $ s_i $ do not use.
In that case, the resulting state will always be unique,
and not in the intersection. This problem can be solved by
erasing temporary variables. In fact, this is the reason
that $ \staterase $ was created.

We do not store any scope
information for local variables in statements.
This implies that local variables
in principle remain visible, once initialized.
This may lead to
space inefficiency when statements are naively translated into $ C $ or
$ C^{++}, $ which allow more restricted scopes. This
problem can be solved by
doing a stack based liveness analysis before translation.
\begin{definition}
   \label{Def_live}
   We define $ {\bf live}_{\bf in}(\pi) $ and
   $ {\bf live}_{\bf out}(\pi) $ as the least
   fixed points of the propagation rules below. The intuitive
   meaning is: $ {\bf live}(\pi)_{\bf in/out} $ are
   the variables whose value may be later used in a computation
   that passes through $ \pi. $
   \begin{itemize}
   \item
      If $ \pi $ contains a statement of form
      $ v \leftarrow t, $ then let $ V $ be the
      set of variables occurring in flatterm $ t. $
      \[ \Phi_{\bf in}(\pi) =
         ( \Phi_{\bf out}(\pi) \backslash \{ v \} ) \cup V. \]
   \item
      If $ \pi $ contains a statement of form
      $ \statisfalse \, v, \ \statistrue \, v, \ \statreturn \, v $
      then
      \[ {\bf live}_{\bf in}( \pi ) =
                {\bf live}_{\bf out}( \pi ) \cup \{ v \}. \]

   \item
      If $ \pi $ contains a statement of form
      $ v. {\rm push\_back}( w ), $ then
      \[ {\bf live}_{\bf in}(\pi) = {\bf live}_{\bf out}(\pi) \cup
              \{ w, v \}. \]
   \item
      If $ \pi $ contains $ \statnop, $ then
      $ {\bf live}_{\bf out}( \pi ) \subseteq {\bf live}_{\bf in} ( \pi ). $
   \item
      If $ \pi $ contains $ \statbranch( s_1, \ldots, s_n ), $
      then let $ \pi_1, \ldots, \pi_n $ be the positions of
      $ s_1, \ldots, s_n. $
      For every $ i, $ we have
      $ \Phi_{\bf in}(\pi_i) \subseteq \Phi_{\bf in}( \pi ), $ and
      $ \Phi_{\bf out}( \pi ) \subseteq \Phi_{\bf out}(\pi_i). $

   \end{itemize}
  A variable $ v $ is \emph{live} at some program point $ \pi $
   if $ v \in {\bf live}_{\bf in}(\pi). $
\end{definition}
Instead of defining $ {\bf live}(\pi) $ as a set of variables, one can also
define $ {\bf use}(\pi) $ as a set of pairs $ ( v, \pi'), $
indicating that $ \pi' $ is a point where the value of $ v $ is potentially used.
The definition is only slightly more complicated than Definition~\ref{Def_live} above.

\section{Range States}
\label{Sect_range_states}

We will map the values of variables into their states.
Arrays will be mapped into states representing a multiset. 
We want to be able to handle in-place array conversions, where
one goes through an array with an index
and changes the properties of the elements.
In order to make this possible, it is not sufficient to 
keep track of single variables, because for a given index $ i, $ 
the values on positions before $ i $ may have different properties,
than the values on positions after $ i. $ 
We will therefore map pairs $ ( v, i ), $ where $ i $ is a potential
index of $ v $ separately. We will map $ (v,i) $ into a triple
of states, where the first element is the multiset of values before $ i, $ 
the second element represents the multiset of values on $ i, $
and the third element represents the multiset of values after $ i. $ 
Since we are possibly working on $ v[i], $ its state has to be stored
separately.

We will keep track of the
relative order between variables that are a potential index,
as well as $ 0 $ and the sizes of arrays.
The simplest ways to do this is by keeping a matrix of potential
differences.
If two indices are equal, their difference is zero.
If the first index is bigger than the second, the difference is positive.
\begin{definition}
   \label{Def_diff_set}
   A \emph{difference} has one of the following forms:
   $ -2^{*}, -1, 0, 1, 2^{*}. $
   A \emph{difference set} is a set of differences.
   We write $ \top $ for the full difference set
   $ \{ -2^{*}, -1, 0, 1, 2^{*} \}. $
\end{definition}
The intended meaning of $ -2^{*} $ is '-2 or less'.
The intended meaning of $ 2^{*} $ is '2 or more'.
Difference sets have two levels of uncertainty:
The first is in the semantics of $ -2^{*} $ and $ 2^{*}, $
which represent a set of possible difference values.
The second is the the fact that a set may contain more than one
element. This ugly asymmetry could be removed by replacing
$ -2^{*} $ by $ -2,-3,-4, \ldots, $ but we want the difference
sets to be finitely representable. We don't know a better
alternative to Definition~\ref{Def_diff_set}.

\begin{definition}
   We define two operations on difference sets:
   The first is \emph{negation}:
   If $ D $ is a difference set, then
   $ -D = \{ -d \mid d \in D \}, $
   using $ - $ on differences as defined in Figure~\ref{Fig_negation_oplus}.

   If $ D_1 $ and $ D_2 $ are difference sets, then
   $ D_1 + D_2 $ is defined as
   $ \bigcup \{ \, d_1 \oplus d_2 \mid
                   d_1 \in D_1 \mbox{ and } d_2 \in D_2 \, \}, $
   where $ \oplus $ is defined in Figure~\ref{Fig_negation_oplus}.
   Note that the entries in the table are sets, so that it is possible
   to take their union.
\end{definition}

\begin{figure}
   \caption{Unary $-$ and binary $ \oplus $}
   \label{Fig_negation_oplus}
   \[
      \begin{array}{l||c|c|c|c|c|}
         - & -2^{*} & -1 & 0 & 1 & 2^{*} \\
         \hline
         \hline
         & 2^{*} & 1 & 0 & -1 & -2^{*} \\
         \hline
      \end{array}
      \hspace*{1.0cm}
      \begin{array}{l||c|c|c|c|c|}
         \oplus & -2^{*} & -1 & 0 & 1 & 2^{*} \\
         \hline
         \hline
         -2^{*} & -2^{*} & -2^{*} & -2^{*} & -2^{*},-1 & \top \\
         \hline
         -1 & -2^{*} & -2^{*} & -1 & 0 & 1, 2^{*} \\
         \hline
         0 & -2^{*} & -1 & 0 & 1 & 2^{*} \\
         \hline
         1 & -2^{*}, -1 & 0 & 1 & 2^{*} & 2^{*} \\
         \hline
         2^{*} & \top & 1, 2^{*} & 2^{*} & 2^{*} & 2^{*} \\
         \hline
      \end{array}
   \]
   The entries in the table are always sets. Remember that $ \top $ denotes
   the complete set $ \{ -2^{*}, -1, 0, 1, 2^{*} \}. $
\end{figure}
\begin{definition}
   Let $ I $ be a finite set.
   A \emph{difference matrix} $ \Delta $ \emph{over} $ I $
   is a mapping of $ I \times I $ to
   difference sets.
   We write $ \Delta.{\bf at}(i_1,i_2) $ for
   $ \Delta(i_1,i_2). $
\end{definition}

Possible indices are variables that can have type $ {\bf u64}, $
as well as $ 0 $ and expressions of $ {\bf len}(s), $ where
$ s $ is a variable that holds an array term.
Note that the elements of $ I $ are symbolic: $ I $ consists
of expressions, not of their values.

We will now define a few operations on difference matrices.
We define them imperatively, instead of declaratively.
In other words, we define them as operations that modify
an existing difference matrix, instead of functions
that create a new one.
The reason for this choice is the following:
Below Definition~\ref{Def_range_state},
we will define the effect of a statement on a range state.
Due to overapproximation, and
the essential non-deterministic nature of high-level computation,
the effect of a statemenet is non-deterministic.
After having tried several
ways to specify the possible successor states, we concluded that
a description using non-deterministic computation is the simplest.
In an earlier version, we gave a recursive definition of the 
effect relation $ \vdash $ by means of Horn clauses,
but it turned out rather unreadable. Using imperative language
is somewhat natural, because in the end we are modeling an imperative
language.  

\begin{definition}
        \label{Def_operations_on_difference_matrix}
   Let $ \Delta $ be a difference matrix over an index set $ I. $
   \begin{itemize}
   \item
      $ \Delta. {\bf contains}(i) $ returns true if
      $ \Delta $ has a difference set for $ i. $ 
   \item
      $ \Delta. {\bf at}(i) $ returns the difference
      set for $ i. $ 
   \item
      If $ i \not \in I, $ then
      $ \Delta. {\bf extend}(i) $ modifies $ \Delta $ into
      a difference matrix $ \Delta' $ over $ I \cup \{ i \}, $ s.t.
      whenever both $ i_1, i_2 \not = i, $ then $ \Delta'. {\bf at}(i_1,i_2) =
      \Delta. {\bf at}( i_1,i_2). $
      If either of $ i_1, i_2 = i, $ then $ \Delta'.{\bf at}(i_1,i_2) = \top. $
   \item
      $ \Delta. {\bf assign}( i_1, i_2, d ) $ changes $ \Delta $
      into a difference matrix $ \Delta', $ s.t.
      $ \Delta'.{\bf at}(i_1,i_2) = d, $ and
      $ \Delta'.{\bf at}(i'_1,i'_2) = \Delta.{\bf at}(i'_1,i'_2), $
      whenever $ i'_1 \not = i_1 $ or $ i'_2 \not = i_2. $
   \item
      If $ i \in I, $ then
      $ \Delta. {\bf erase}(i) $ is obtained from $ \Delta $
      by making $ \Delta.{\bf at}(i,i_2) $ and $ \Delta.{\bf at}(i_1,i) $
      undefined, for all $ i_1,i_2 \in I. $ 
   \item
      $ \Delta. {\bf isconsistent}( ) $ returns true if
      there are no $ i_1, i_2 \in I, $ s.t.
      $ \Delta. {\bf at}(i_1,i_2) = \emptyset. $
   \end{itemize}
\end{definition}

\begin{definition}
   Let $ \Delta $ be a difference matrix. We define
   a procedure $ {\bf close}( \Delta ) $ that
   closes $ \Delta $ under reflexivity, transitivity and antisymmetry:
   $ {\bf close}( \Delta ) $
   is obtained by applying the following
   operations, until no further change is obtained:
   \begin{itemize}
   \item
      For every $ i, $ s.t. $ \Delta. {\bf contains}(i): $ \\
      \hspace*{1.0cm} $ 
         \begin{array}{l}
            \Delta. {\bf assign}(i,i, \{ 0 \} ); \hspace*{1.0cm}
            \Delta. {\bf assign}( 0, i, \,
               \Delta. {\bf at}(0,i) \cap \{ 0,1,2^{*} \} ); \\ 
            \Delta. {\bf assign}( i, {\bf len}(v),
               \Delta. {\bf at}( i, {\bf len}(v)) \cap \{ 0,1,2^{*} \} \, ) \\
         \end{array}
      $ 
   \item
      For every $ i_1,i_2, $ s.t. $ \Delta. {\bf contains}(i_1) $ and
      $ \Delta. {\bf contains}(i_2): $ \\
      \hspace*{1.0cm}
      $ 
           \Delta. {\bf assign}(i_1,i_2, \,
           \Delta. {\bf at}(i_1,i_2) \cap -\Delta. {\bf at}(i_2,i_1) \, ). 
      $
   \item
      For every $ i_1,i_2,i_3, $ s.t. $ \Delta. {\bf contains}(i_1), \
      \Delta. {\bf contains}(i_2), $ and $ \Delta. {\bf contains}(i_3): $ \\
      \hspace*{1.0cm} 
      $ \Delta.{\bf assign}(i_1,i_3), \
           \Delta. {\bf at}(i_1,i_3) \cap (
           \Delta. {\bf at}(i_1,i_2) + \Delta. {\bf at}(i_2,i_3)) \ ). $ 
   \end{itemize}
\end{definition}
We now define state sequences and a couple of operations
on them.
As for difference matrixes, we define the operations imperatively.
The operations on state sequences are not yet non-deterministic,
but they can fail.
\begin{definition}
   \label{Def_state_seq}
   Let $ {\cal A} = ( Q, \delta, N ) $ be a DMA.
   A \emph{state sequence} over $ {\cal A} $ is a
   (possibly empty) sequence of states of $ {\cal A}. $
   We define the following operations on state sequences:
   \begin{itemize}
   \item
      $ \overline{q}. {\bf push\_back}(q') $ replaces
      $ \overline{q} = (q_1, \ldots, q_n ) $ by
      $ (q_1, \ldots, q_n, q' ), $ and
      $ \overline{q}. {\bf push\_front}(q') $ replaces
      $ \overline{q} = (q_1, \ldots, q_n) $ by
      $ ( q', q_1, \ldots, q_n ). $
   \item
      If $ \overline{q} $ has form $ ( q_1, \ldots, q_n ) $ with $ n > 0, $
      then $ \overline{q}. {\bf pop\_back}( ) $ replaces
      $ \overline{q} $ by $ ( q_1, \ldots, q_{n-1} ), $
      and $ \overline{q}. {\bf pop\_front}( ) $ replaces
      $ \overline{q} $ by $ ( q_2, \ldots, q_n ). $
   \item
      If $ \overline{q} $ has form $ ( q_1, \ldots, q_n ) $ with $ n > 0, $
      then $ \overline{q}. {\bf back}( ) $ equals $ q_n, $
      and $ \overline{q}. {\bf front}( ) $ equals $ q_1. $
   \end{itemize}
\end{definition}
The purpose of a state sequence is to approximate the
data terms at positions $ v[i], v[i+1], \ldots. $ 
In fact, we currently do not use $ v[i+1], \ldots $ and looking at
$ v[i] $ is optional. Technically, we could use a 
single optional value. We defined Definition~\ref{Def_state_seq} a bit
more general, than needed. 

\begin{definition}
        \label{Def_operations_on_DMA_states}
   Let $ {\cal A} = ( Q, \delta, N ) $ be a DMA.
   \begin{itemize}
   \item
      Let $ q $ be a state of $ {\cal A}. $
      Operation $ q. {\bf insert}(q') $ replaces
      $ q $ by $ \delta( q;q' ). $
   \item
      Method $ q. {\bf extract}( ) $ non-deterministically
      does the following:
      Find $ q', q_1, $ s.t. $ \delta(q';q_1) = q. $
      Replace $ q $ by $ q' $ and return $ q_1. $
   \end{itemize}
\end{definition}
The second operation is one of the sources of non-determinism
in the effect of an operation, the other is uncertainty about the
relative order of indices.
\begin{definition}
   \label{Def_range_state}
   Let $ {\cal A} $
   be a an abstraction automaton.
   An \emph{range state} is a pair of form $ ( M, \Delta ), $
   s.t.
   \begin{itemize}
   \item
      $ M $ is a finite partial function, which
      \begin{itemize}
      \item
         maps single variables to $ Q. $
      \item
         maps pairs $ (v,i) $ consisting of two variables $ v,i $ to
         triples $ ( q_{1}, \overline{q}, q_{2} ). $
      \end{itemize}
      Both components can be partial.
   \item
      $ \Delta $ is a difference matrix over a set
      $ I $ consisting of variables, $ 0, $
      and expressions of form $ {\bf len}(v), $
      with $ v $ a variable.
   \end{itemize}
   We define the following methods of $ M: $
   \begin{itemize}
   \item
      $ M. {\bf contains}(v), \ M. {\bf contains}(v,i). $
      True if $ M $ contains a value for variable $ v, $
      or for the pair $ (v,i). $
   \item
      $ M. {\bf at}(v), \ M. {\bf at}(v,i). $ Returns the value
      for the variable $ v $ or the pair $ (v,i) $ in $ M. $ 
   \item
      $ M. {\bf assign}(v,d), \ M. {\bf assign}(v,i,d). $ Sets the value
      of variable $ v, $ or the pair $ (v,i) $ to $ d. $
   \item
      $ M. {\bf erase}(v), \ M. {\bf erase}(v,i). $ Remove the
      value for variable $ v, $ or the pair $ (v,d). $
   \end{itemize}
\end{definition}
The intended meaning of $ q_{1}, \overline{q}, q_{2} $
is that they represent the multisets of values
in $ v[ 0 \cdots i \rangle $, $ v[i] $, and $ v \langle i \cdots \rangle $, respectively.

For each program point, we will construct a set of
range states that are possible at that point,
using Definition~\ref{Def_semantics}.
As usual, it cannot be avoided that there will be abstract states
attached to a program point that are in fact unreachable.

\noindent
For simplicity, we will always assume that all assignments to
variables (not to a member) are initializations.
We explain the relation between assignment and
comparison. Since our computational model is non-deterministic,
assignment can be viewed as a special case of comparison,
where nothing is known about the assigned variable.
More precisely, the assignment $ i \leftarrow 0 $ can be viewed
as two statements: {\bf (1)} Initialize $ i $ with a random value.
{\bf (2)} Verify that $ i $ equals $ 0, $ failing if it is not.
Since computation is non-deterministic, $ i $ will be equal
to $ 0 $ in a succesful computation.
It is convenient to treat assignment as 
comparison as much as possible, and to deal with
known equalities in a postprocessing state. This has the advantage
that the common parts of assignment and comparison can be treated
by a single procedure.

We list possible statements, and their effect on range states.
Due to lack of space, the list is incomplete. 
\begin{itemize}
\item
   The statement $ i \leftarrow 0 $ has the following effect: \\
   \hspace*{1.0cm}
   $ 
      \Delta. {\bf extend}(i); \hspace*{0.5cm} 
      \Delta. {\bf assign}(i,0, \{ 0 \} ); \hspace*{0.5cm}
      {\bf close}( \Delta ); \hspace*{0.5cm}
      M. {\bf assign}( i, Q_{\cal A}(0) );
   $ 
\item
   A statement of form $ i \leftarrow {\bf len}(v) $
   has the following effect: \\
   \hspace*{1.0cm} 
   $ 
         \Delta. {\bf extend}(i); \hspace*{0.5cm}
         \Delta. {\bf assign}(i, {\bf len}(v), \{ 0 \} ); \hspace*{0.5cm}
         {\bf close}( \Delta );
   $
\item
   A statement of form $ i \leftarrow j $ 
   has the following effect: \\
   \hspace*{1.0cm} 
   $
         \Delta. {\bf extend}(i); \hspace*{0.5cm}
         \Delta. {\bf assign}(i, j, \{ 0 \} ); \hspace*{0.5cm}
         {\bf close}( \Delta ); 
   $
\item
   A statement of form $ w \leftarrow v[i] $
   has the following effect: \\
   \hspace*{1.0cm} 
   $ \begin{array}{l}
      {\rm let} \ ( q_1, \overline{q}, q_2 ) = M. {\bf at}(v,i) \\
      \mbox{if} \ \| \overline{q} \| = 0 \ \mbox{then} \\
      \hspace*{0.5cm} \overline{q}. {\bf push\_back}( q_2. {\bf extract}( )) \\
      M. {\bf assign}( w, \overline{q}. {\bf front}( )) \\
      M. {\bf assign}( v,i, ( q_1, \overline{q}, q_2 ) ) \\
      \end{array}
   $ \\ \\ 
   If the triple at $ M. {\bf at}(v,i) $ has a value for
   $ v[i], $ we can use it. Otherwise, we create it first. 
\item
   A statement of form $ j \leftarrow i + 1 $ has the following effect:
   \[
      \begin{array}{l}
      \Delta. {\bf assign}( j,i, \{ 1 \} ); \ \ {\bf close}( \Delta ) \\
      {\rm let} \ ( q_1, \overline{q}, q_2 ) = M. {\bf at}(v,i) \\
      \mbox{if} \ \| \overline{q} \| = 0 \ \mbox{then} \\
      \hspace*{0.5cm} \overline{q}. {\bf push\_back}( q_2. {\bf extract}( )) \\
      M. {\bf assign}( v,j, ( \delta( q_1; \overline{q}. {\bf front}()), q_2 )) \\ 
      \end{array}
   \] 
   We take $ M.{\bf at}(v,i), $ create the value at $ v[i] $ if it is not
   present, and transfer it to $ v[ 0 \cdots j \rangle. $ 
   Note that we are taking a slight shortcut.
   The problem is that $ i + 1 $ is not a flatterm.
   In reality, one must check for assignments of form
   $ j \leftarrow i_1 + i_2, $ where either
   $ \Delta.{\bf at}(i_1,0) = \{ 1 \}, $ or
   $ \Delta.{\bf at}(i_2,0) = \{ 1 \}. $ 
\item
   A statement of form $ v[i] \leftarrow w $ has the following effect:
   For every variable $ j $ that is a potential index of $ v, $ do: 
   \[ 
      \begin{array}{l} 
         {\rm require} \ \Delta(i, {\bf len}(v)) \cap \{ 0,1,2^{*} \} \not = \emptyset \\
         {\rm let} \ ( q_1, \overline{q}, q_2 ) = M. {\bf at}(v,j) \\
         {\rm nondeterministic} \\ 
         \hspace*{0.5cm} \bullet \ 
              {\rm require} \ \Delta.{\bf at}( i,j ) \cap \{ -1, -2^{*} \} \not = \emptyset \\
         \hspace*{0.8cm} q_1. {\bf extract}( ); \,
                         q_1. {\bf insert}( Q_{\cal A}( M. {\bf at}(w))); \\
         \hspace*{0.5cm} \bullet \ {\rm require} \
             0 \in \Delta.{\bf at}(i,j) \mbox{ and } 
                 \| \overline{q} \| = 1 \\
         \hspace*{0.8cm} {\rm let} \ \overline{q} = Q_{\cal A}( M. {\bf at}(w)) \\
         \hspace*{0.5cm} \bullet \  
          {\rm require} \ \Delta(i,j) \cap \{ 1,2^{*} \} \not = \emptyset 
                      \mbox{ or }
                          0 \in \Delta( i, j ) \mbox{ and } 
                          \| \overline{q} \| = 0 \\
          \hspace*{0.8cm} q_2. {\bf extract}( ); \,
                          q_2. {\bf insert}( Q_{\cal A}( M. {\bf at}(w))) \\
          M. {\bf assign}( v, i, \, ( q_1, \overline{q}, q_2 ) \, ) \\ 
       \end{array}
   \] 
   A require statement fails quietly if its condition is not met.
   The nondetermistic statement branches into three options 
   each of which can succeed or fail independently.
   The code may be surprising at first. First observe that $ i $ 
   will be included in the for loop. For $ i, $ the 
   middle option is entered if $ \overline{q} $ is non-empty,
   and $ \overline{q} $ will be replaced by 
   $ Q_{\cal A}( M. {\bf at}(w)). $ If $ \overline{q} $ is empty,
   the third option will be entered, and a random element from
   $ q_2 $ wil be replaced by $ Q_{\cal A}( M. {\bf at}(w)). $ 
   The assignment to $ v[i] $ may also affect 
   $ v[ 0, \ldots j \rangle $ or $ v[ j \ldots \rangle \rangle $ 
   for other variables $ j, $ for which the interval
   $ [ 0, \ldots j \rangle $ or $ [ j, \ldots, \rangle $ might
   contain $ i. $ 
\item
   A statement of form $ v \leftarrow w.f $ has the following effect:
   Let 
   $ i $ be the index of field $ f. $ 
   Let $ q = Q_{\cal A}(w). $  
   Non-deterministically pick a tuple $ ( q_1, \ldots, q_n ), $ for
   which $ \delta( q_1, \ldots, q_n ) = q. $
   $ M. {\bf assign}(v, q_i ). $ 
\item
   A statement of form $ v \leftarrow {\bf lt}(w_1,w_2) $
   has the following effect:
   \[
       \begin{array}{l}
          {\rm nondeterministic} \\
             \hspace*{0.5cm} \bullet \ \Delta. {\bf assign}( w_1, w_2,
               \Delta. {\bf at}( w_1, w_2 ) \cap \{ -2^{*}, -1 \} \, ); \ 
            {\bf close}( \Delta ) \\
            \hspace*{0.8cm} \mbox{require} \ \Delta. {\bf isconsistent}( )) \\
            \hspace*{0.8cm} M. {\bf assign}(v, Q_{\cal A}( {\bf t} )) \\
            \hspace*{0.5cm} \bullet \ \Delta. {\bf assign}( w_1, w_2,
               \Delta. {\bf at}( w_1, w_2 ) \cap \{ 0,1,2^{*} \} \, ); \ 
            {\bf close}( \Delta ) \\
            \hspace*{0.8cm} \mbox{require} \ \Delta. {\bf isconsistent}( )) \\
            \hspace*{0.8cm} M. {\bf assign}( v, Q_{\cal A}( {\bf f} )) \\
         \end{array}
      \]
      The other comparison operators are analogeous.
      Executing the code above results in either one or two successor states,
      dependent on whether assuming $ w_1 < w_2 $ or $ w_1 \geq w_2 $ is consistent
      with $ \Delta. $
\item
   A statement of form $ v \leftarrow F(w), $ with
   $ F $ a unary state condition has the following effect: \\
   \hspace*{1.0cm}
   $ \begin{array}{l}
         \mbox{if } F( Q_{\cal A}(w)) \\
         \hspace*{0.5cm} M. {\bf assign}(v, Q_{\cal A}({\bf t})) \\
         \mbox{else} \\
         \hspace*{0.5cm} M. {\bf assign}(v, Q_{\cal A}({\bf f})) \\
      \end{array}
   $ 
\end{itemize}
It remains to define a postprocessing procedure that must be called
after every operation on a range state.
Its purpose is to check for equalities inside the
difference matrix $ \Delta, $ and apply them. 
Similar to the description of the effect of a statement, we
describe the procedure in imperative fashion. 
\begin{definition}
   \label{Def_postproc}
   First define a function $ {\bf part}( q_1, \overline{q}, q_2 ) $ as
   follows: \\ 
   \hspace*{0.5cm}
   $ 
         \begin{array}{l}
            q = q_2; \ 
            \mbox{for} \ q' \in \overline{q} \mbox{ do} \\
            \hspace*{0.5cm}
               q = \delta( q'; q ) \\
            {\rm return} \ ( q_1, q ) 
         \end{array}
   $ \\ 
   The purpose of function $ {\bf part} $ is to partition
   the triple $ (q_1, \overline{q}, q_2) $ by pushing the middle component
   $ \overline{q} $ back into $ q_2. $ We use {\bf part} when
   comparing different $ M. {\bf at}(v,i), $ where we want
   to ignore the middle components.
   \begin{itemize}
   \item
      For every variable $ i, $ s.t.
      $ \Delta. {\bf at}(i,0) = \{ 0 \}, $ for every variable $ v, $
      s.t. $ M. {\bf contains}(v), $ do the following: 
      If $ M.{\bf contains}(v,i), $ then 
      require that $ {\bf part}( M. {\bf at}(v,i)) $ is equal to
      $ ( \delta( ), M. {\bf at}(v)). $ 
      Otherwise, $ M. {\bf assign}(v,i, ( \delta( ), ( ), M. {\bf at}(v)). $ 
   \item
      For every pair of variables $ i $ and $ v, $ s.t 
      $ \Delta. {\bf at}(i, {\bf end}(v)) = \{ 0 \} $ do the following:
      If $ M. {\bf contains}(v,i), $ then require that
      $ {\bf part}( M. {\bf at}(v,i)) $ is equal to
      $ ( M. {\bf at}(i), \delta()). $
      Otherwise,  
      $ M. {\bf assign}( v,i, ( M. {\bf at}(v), ( \, ), \delta( )). $ 
   \item
      For every pair of variables $ i, j, $ s.t.
      $ \Delta.{\bf at}(i,j) = \{ 0 \}, $ for every
      $ v, $ s.t. $ M. {\bf contains}(v), $ do the following:

      If $ M. {\bf contains}(v,i) $ holds, but not
         $ M. {\bf contains}(v,j), $ then 
      $ M. {\bf assign}(v,j, M. {\bf at}(v,i)). $ 
      Similarly, if
      $ M. {\bf contains}(v,j) $ holds, but
      $ M. {\bf contains}(v,i) $ does not hold, then
      $ M. {\bf assign}(v,i, M. {\bf at}(v,j)). $
      Otherwise, require that 
      $ {\bf part}( M. {\bf at}(v,i)) $ equals 
      $ {\bf part}( M. {\bf at}(v,j)). $ 
   \item
      For every variable $ v, $ s.t. $ M. {\bf contains}(v), $ 
      but $ v $ does not occur in $ {\bf live}_{\bf out}( \pi ), $
      call $ M. {\bf erase}(v). $
      For every pair of variables $ v,i, $ s.t.
      $ M. {\bf contains}(v,i), $ but
      either $ v $ or $ i $ does not occur in $ {\bf live}_{\bf out}(\pi), $
      call $ M. {\bf erase}(v,i). $
   \end{itemize}
\end{definition}
The description of the effects of statements that we gave is necessarily
incomplete, due to lack of space. In addition, there are many small 
variations possible. What we presented should be viewed as a toolbox
from which type checking procedures can be obtained.

We did not explain the type checking procedure itself. 
We give a short, intuitive overview. 
Initially, for every ambiguous statement, all candidate overloads
are inserted into the statement, joined by a {\bf resolve} statement.
We use Definition~\ref{Def_semantics} combined
with the range state transformations in Section~\ref{Sect_range_states}
to obtain a fixpoint containing combinations of possible states. 
After that, we cut out non-viable overloads 
from the flow graph and repeat the procedure until we reach
a fixed point where all overloads are viable. 
At this point, it must be the case that every
resolve statement is non-empty, and contain one candidate 
that is strictly more specific than all of the other candidates.
This can be checked with the procedure sketced below 
Definition~\ref{Def_acc_proj}.

\section{Conclusions and Future Work}
\label{Sect_concl_fut}

We presented a context-aware type system and a type-checking procedure
that could be used for implementing an experimental imperative
programming language working on tree-like data structures
which can have different forms based on the way they are constructed.
The type system can determine forms of data terms from a program
context and ensure the correctness of field accesses and function applications.
We expect that, using context-aware typing, many casts that 
are currently present in statically typed,
imperative programs, can be removed. 
In \cite{AlexanderAikenBrianMurphy1991_1} a related approach for
functional programs is considered, 
but we consider imperative programs, and handle arrays.

At this moment, we do not have an implementation of the
type checker. We have no empirical data about its feasibility.
Our first priority is develop a working prototype. 
This would allow practical validation of the system and
provide insights into real-world challenges.
In particular, it will be important to assess the scalability
and performance of the system with real-world examples.
At this moment, we are not trying to prove theoretical upperbounds
to the checking procedure, like for example by 
\cite{AlexanderAikenDexterKozenMosheVardi1994} or \cite{ChaPo2002},
because such upperbounds tend to be
fragile, which implies that small changes in the system can cause
dramatic differences in complexity. Once we have a stable system,
we will look into theoretical complexity.

We did not discuss how tree automata
are obtained from real type conditions. In particular, we did not
discuss one important practical problem, namely how to treat data 
trees that are not part of a precondition. The automaton must
be set up in such a way that some mininal information about the
types contained in it, need to be preserved. This can be done by
adding so called \emph{overflow states}. Different heuristics
for adding overflow states are possible. It has to be tried by
experimenting which heuristic gives the best results.

At present, the sets $ \Phi_{\bf in} $ and $ \Phi_{\bf out} $ are 
constructed as sets of independent states. Since a single state can
be represented as a conjunction of atoms of form
$ \sigma(v) = t, $ a set of possible states can be
viewed as a single formula in disjunctive normal form.
Using conjunctive normal form would result in a much
more compact representation, but in the process, information
about dependencies between variables will be lost. In most cases,
this is unproblematic, but it matters in some cases, e.g.
after one has checked that two selector variables are equal,
or when one has stored the result of a comparison in a boolean variable.
Future work needs to focus on finding a compact representation that
keeps just enough information about dependencies.
Perhaps one could use a decision tree based on a few binary predicates.

\section{Acknowledgements}

This work gained from discussions with Witold Charatonik.
We thank Nazarbayev University for
supporting this research through the Faculty Development Competitive
Research Grant Program (FDCRGP) number 021220FD1651.

\bibliographystyle{plain}
\bibliography{types,../../administration/cv/pub.bib}

\end{document}
